\chapter[Paraboliczne równania ewolucyjne]{Paraboliczne równania ewolucyjne: rozwiązania słabe, metoda Galerkina i regularność}
\label{chap:pde20}

\section{Od równania ciepła do współczynników mierzalnych}
W rozdziale~4 rozwiązywaliśmy równanie ciepła w przestrzeni całej
przy użyciu jądra Gaussa. Obecnie rozważymy znacznie ogólniejsze
zagadnienie na ograniczonym obszarze $\Omega\subset\R^d$:
\begin{equation}\label{eq:pde20-eq}
 \partial_t u-\operatorname{div}(A(t,x)\nabla u)=f
 \quad\text{w }(0,T)\times\Omega,\qquad
 u|_{\partial\Omega}=0,\qquad u(0)=u_0.
\end{equation}
Macierz $A$ może być tylko mierzalna zarówno w czasie, jak i w przestrzeni.
W takich warunkach nie należy oczekiwać klasycznych pochodnych drugiego
rzędu. W przeciwieństwie do zagadnienia eliptycznego rozwiązanie zależy
od czasu, ale nie zakładamy istnienia pochodnej $u_t$ o wartościach w $L^2(\Omega)$.

\begin{definicja}[Dane i eliptyczność paraboliczna]\label{def:pde20-data}
Niech $\Omega\subset\R^d$ będzie ograniczonym obszarem Lipschitzowskim,
$0<T<\infty$, $H=L^2(\Omega)$, $V=H_0^1(\Omega)$, $V^*=H^{-1}(\Omega)$.
W całym rozdziale rozpatrujemy rzeczywiste przestrzenie
Hilberta, a $V^*$ oznacza przestrzeń ciągłych
funkcjonałów liniowych na $V$.
Rozpatrujemy rzeczywistą macierz $A=(a_{ij})\in
L^\infty((0,T)\times\Omega;\R^{d\times d})$ taką, że dla pewnych
$0<\lambda\le\Lambda<\infty$, dla prawie wszystkich $(t,x)$ i każdego
$\xi\in\R^d$ zachodzi
\[
 \lambda|\xi|^2\le A(t,x)\xi\cdot\xi,
 \qquad |A(t,x)\xi|\le\Lambda|\xi|.
\]
Nie wymagamy symetrii macierzy $A$. Dla $v,w\in V$ określamy
$a(t;v,w)=\int_\Omega A(t,x)\nabla v(x)\cdot\nabla w(x)
\,\mathrm d\lambda_d(x)$. Przyjmujemy $u_0\in H$ oraz
$f\in L^2(0,T;V^*)$; dualność $V^*,V$ oznaczamy przez
$\langle\cdot,\cdot\rangle$.
\end{definicja}


\begin{center}
\begin{tikzpicture}[scale=0.9,>=Latex]
  \fill[blue!4] (0,0) rectangle (6,3.4);
  \draw[very thick,blue!45!black] (0,0) -- (0,3.4);
  \draw[very thick,blue!45!black] (6,3.4) -- (6,0);
  \draw[dashed,gray] (0,3.4) -- (6,3.4);
  \draw[very thick,teal!65!black] (0,0)--(6,0);
  \draw[->] (-0.4,0)--(6.8,0) node[right] {$x$};
  \draw[->] (0,-0.25)--(0,4.1) node[above] {$t$};
  \node[below] at (3,-0.07) {$u(0,x)=u_0(x)$};
  \node[rotate=90,above] at (0,1.8) {$u(t,0)=0$};
  \node[rotate=90,below] at (6,1.8) {$u(t,1)=0$};
  \node[above] at (3,3.4) {$t=T$ (wartość nie jest zadawana)};
  \node at (3,1.75) {$Q_T=(0,T)\times(0,1)$};
  \node[below] at (0,-0.4) {$0$};
  \node[below] at (6,-0.4) {$1$};
\end{tikzpicture}

\small Rysunek: jednowymiarowy przykład ($\Omega=(0,1)$). Podstawa i brzeg boczny walca parabolicznego.
Dane początkowe i warunek Dirichleta określa się na różnych częściach brzegu.
\end{center}

Normę w $V$ ustalamy jako $\|v\|_V=\|\nabla v\|_{L^2(\Omega)}$.
Dzięki nierówności Poincarégo (por. sekcję~\ref{sec:pde18-intro})
jest ona równoważna zwykłej normie $H^1$. Oznaczamy przez
$C_P$ stałą w nierówności $\|v\|_H\le C_P\|v\|_V$.
Mamy trójkę Gelfanda
\[
 V\hookrightarrow H\cong H^*\hookrightarrow V^*,
 \qquad h\longmapsto[v\mapsto(h,v)_H].
\]
Pierwsze zanurzenie jest ciągłe i ma gęsty obraz,
bo zawiera $C_c^\infty(\Omega)$. Drugie jest iniekcją
dzięki tej gęstości, a nierówność Cauchy'ego--Schwarza
daje $\|h\|_{V^*}\le C_P\|h\|_H$.
Jego obraz jest gęsty w $V^*$: gdyby nie był, istniałby
niezerowy ciągły funkcjonał na $V^*$ zerujący się na $H$.
Z refleksywności przestrzeni Hilberta $V$ taki funkcjonał
miałby postać $g\mapsto\langle g,v\rangle$ dla pewnego
$v\in V$. Warunek $(h,v)_H=0$ dla wszystkich $h\in H$
dawałby jednak $v=0$.

Dla prawie każdego $t$ forma $a(t;\cdot,\cdot)$ spełnia
\[
 |a(t;v,w)|\le\Lambda\|v\|_V\|w\|_V,\qquad
 a(t;v,v)\ge\lambda\|v\|_V^2.
\]
Pierwsze oszacowanie wynika z nierówności
Cauchy'ego--Schwarza i ograniczenia macierzy, drugie
z eliptyczności. Niech $\mathcal A(t):V\to V^*$ będzie
operatorem określonym przez
$\langle\mathcal A(t)v,w\rangle=a(t;v,w)$.
Dla ustalonych $v,w\in V$ mierzalność tej funkcji czasu
wynika z twierdzenia Fubiniego, ponieważ całkowany
iloczyn jest zmajoryzowany przez
$\Lambda|\nabla v||\nabla w|\in L^1(\Omega)$.
Przestrzeń $V$ jest ośrodkową przestrzenią Hilberta
(jako domknięta podprzestrzeń przestrzeni izometrycznie
zanurzonej w $L^2(\Omega)^{d+1}$); również $V^*$
jest ośrodkowa. Mierzalność skalarna $\mathcal A(t)v$
pociąga więc silną mierzalność, na przykład przez
rozwinięcie w przeliczalnej bazie ortonormalnej $V^*$.
Aproksymacja $u(t)$ funkcjami prostymi i jednolite
ograniczenie $\|\mathcal A(t)\|\le\Lambda$ pokazują,
że $\mathcal A(\cdot)u(\cdot)$ jest silnie mierzalna
dla $u\in L^2(0,T;V)$ oraz
$\|\mathcal A(\cdot)u(\cdot)\|_{L^2(0,T;V^*)}
\le\Lambda\|u\|_{L^2(0,T;V)}$. Wszystkie przestrzenie
$L^2(0,T;X)$ rozumiemy w sensie Bochnera.

\begin{uwaga}[Trzy odrębne narzędzia: forma, baza i zwartość]
\label{rem:pde-v43-parabolic-basis}
W dowodzie istnienia wykorzystamy koercywność formy
(znaną z rozdziału~\ref{chap:pde-variational}), bazę własną
pomocniczego $-\Delta_D$ z rozdziału~\ref{chap:pde18}
oraz zwarte zanurzenie z rozdziału~\ref{chap:pde-embeddings}.
Baza laplasjanu nie diagonalizuje na ogół macierzy
$A(t,x)$; służy do zbudowania stabilnych rzutów
$P_n$ w $V$, $H$ i $V^*$. Dlatego ani niesymetryczność,
ani mierzalna zależność $A$ od czasu nie podważają
przedstawionej konstrukcji. Własności samosprzężonej
półgrupy $e^{-tL}$ w dalszej części rozdziału
odnoszą się natomiast do wyraźnie wskazanego przypadku
$L=-\Delta_D$ (lub innego stałego operatora dodatniego
samosprzężonego), nie do ogólnego $A(t,x)$.
\end{uwaga}

\begin{definicja}[Rozwiązanie słabe w przestrzeni energii]
\label{def:pde20-weak}
Rozwiązaniem słabym zagadnienia~\eqref{eq:pde20-eq} nazywamy funkcję
\[
 u\in L^2(0,T;V),\qquad u'\in L^2(0,T;V^*),
\]
której pochodną rozumiemy dystrybucyjnie w czasie:
dla każdej $\eta\in C_c^\infty(0,T)$ zachodzi tożsamość
w $V^*$
\[
 \int_0^T u(t)\eta'(t)\,\mathrm d\lambda_1(t)
 =-\int_0^T u'(t)\eta(t)\,\mathrm d\lambda_1(t).
\]
Tutaj $u$ po lewej stronie utożsamiamy z jego obrazem
w $V^*$. Ponadto rozwiązanie spełnia
$\langle u'(t),v\rangle+a(t;u(t),v)=\langle f(t),v\rangle$
dla każdego $v\in V$ i prawie każdego $t\in(0,T)$, przy czym zbiór
wyjątkowy można wybrać niezależnie od $v$.
Równoważnie mamy $u'(t)+\mathcal A(t)u(t)=f(t)$
w $V^*$ dla prawie każdego $t$. Gdyby równanie
skalarne sformułowano z osobnym zbiorem wyjątkowym
dla każdego $v$, wystarczyłoby wziąć przeliczalny
gęsty podzbiór $V$, zsumować jego zbiory wyjątkowe
i zastosować ciągłość obu stron względem $v$.
Warunek $u(0)=u_0$ dotyczy ciągłego reprezentanta $H$-wartościowego,
którego istnienie wykażemy poniżej.
\end{definicja}

\section{Ciągłość czasowa i tożsamość energii}
\begin{lemat}[Tożsamość energii w trójce Gelfanda]\label{lem:pde20-energy}
Niech $V=H_0^1(\Omega)$, $H=L^2(\Omega)$, a
$u\in L^2(0,T;V)$ ma pochodną $u'\in L^2(0,T;V^*)$.
Wówczas $u$ ma jednoznacznego reprezentanta należącego do
$C([0,T];H)$ i dla wszystkich $0\le s\le t\le T$ zachodzi
\begin{equation}\label{eq:pde20-energy-identity}
 \frac12\|u(t)\|_H^2-\frac12\|u(s)\|_H^2
 =\int_s^t\langle u'(r),u(r)\rangle\,\mathrm dr.
\end{equation}
\end{lemat}
\begin{proof}
Niech $\{e_j\}_{j\ge1}$ będzie ortonormalną bazą $H$ złożoną
z funkcji własnych dodatniego laplasjanu Dirichleta,
$-\Delta_D e_j=\mu_j e_j$, $\mu_j>0$.
Z rozwinięcia formy energii w dowodzie twierdzenia
~\ref{thm:pde18-minmax} wynika, że
$\|v\|_V^2=\sum_{j\ge1}\mu_j|(v,e_j)_H|^2$ dla $v\in V$.
Układ $(\mu_j^{-1/2}e_j)$ jest bazą ortonormalną $V$.
Rzut ortogonalny $P_n$ w $H$ na
$\operatorname{span}\{e_1,\ldots,e_n\}$ jest więc
również rzutem ortogonalnym w $V$.
Na $V^*$ definiujemy zgodne rozszerzenie przez
$\langle P_ng,v\rangle=\langle g,P_nv\rangle$,
czyli $P_ng=\sum_{j=1}^n\langle g,e_j\rangle e_j$
przy utożsamieniu elementów $H$ z funkcjonałami na $V$.
Z twierdzenia Riesza i tożsamości Parsevala w $V$
wynika
\[
 \|g\|_{V^*}^2=\sum_{j\ge1}\mu_j^{-1}
 |\langle g,e_j\rangle|^2.
\]
Zatem $P_n$ jest kontrakcją we wszystkich trzech
przestrzeniach i zbiega silnie do identyczności:
kwadraty norm różnic są ogonami odpowiednich szeregów.
Twierdzenie o zbieżności zmajoryzowanej zastosowane
do kwadratów norm różnic daje
$P_nu\to u$ w $L^2(0,T;V)$ oraz
$P_nu'\to u'$ w $L^2(0,T;V^*)$.

Współczynniki $t\mapsto (u(t),e_j)_H$ należą do $H^1(0,T)$,
ponieważ testując definicję pochodnej czasowej
funkcją $\eta$ i funkcjonałem ewaluacji na $e_j$,
otrzymujemy ich słabe pochodne
$\langle u'(t),e_j\rangle\in L^2(0,T)$.
Wybieramy ich absolutnie ciągłe reprezentanty
(twierdzenie~\ref{thm:pde9-ac}); wtedy $P_nu$
jest ciągłe na $[0,T]$ i jego pochodna to $P_nu'$.
Dla skończonego $n$ zwykła reguła różniczkowania daje
\[
 \|P_nu(t)\|_H^2-\|P_nu(s)\|_H^2
 =2\int_s^t\langle P_nu'(r),P_nu(r)\rangle\,\mathrm dr.
\]
Wykażemy najpierw zbieżność $P_nu$ w $C([0,T];H)$.
Dla $w=P_nu-P_mu$ wybieramy punkt $s\in[0,T]$ taki, że
$\|w(s)\|_H^2\le T^{-1}\|w\|_{L^2(0,T;H)}^2$;
istnieje on, gdyż średnia nie może być mniejsza od wszystkich wartości.
Stosując poprzednią tożsamość na przedziale o końcach $s,t$
i nierówność Cauchy'ego--Schwarza, otrzymujemy
\begin{equation}\label{eq:pde20-cont-est}
 \sup_{0\le t\le T}\|w(t)\|_H^2
 \le T^{-1}\|w\|_{L^2(0,T;H)}^2
 +2\|w'\|_{L^2(0,T;V^*)}\|w\|_{L^2(0,T;V)}.
\end{equation}
Oba składniki po prawej dążą do zera dla $m,n\to\infty$.
Granica jest więc $H$-wartościowo ciągła i równa $u$
prawie wszędzie. Przechodząc do granicy w skończenie wymiarowej
tożsamości energii, korzystamy z jednostajnej zbieżności
wartości końcowych oraz oszacowania
\[
 \begin{split}
 \int_0^T|\langle P_nu',P_nu\rangle-\langle u',u\rangle|
 \,\mathrm d\lambda_1
 &\le \|P_nu'-u'\|_{L^2(0,T;V^*)}\|P_nu\|_{L^2(0,T;V)}\\
 &\quad+\|u'\|_{L^2(0,T;V^*)}\|P_nu-u\|_{L^2(0,T;V)}.
 \end{split}
\]
Prawa strona dąży do zera. To daje~\eqref{eq:pde20-energy-identity}.
Jednoznaczność reprezentanta wynika z ciągłości w $H$:
dwa ciągłe reprezentanty równe prawie wszędzie muszą
być równe w każdym punkcie $[0,T]$.
W szczególności funkcja $t\mapsto\frac12\|u(t)\|_H^2$
jest absolutnie ciągła, a jej pochodną prawie wszędzie
jest $\langle u'(t),u(t)\rangle$, która jest całkowalna
z nierówności Cauchy'ego--Schwarza.
\end{proof}

\section{Istnienie i jednoznaczność: pełna konstrukcja Galerkina}
\begin{twierdzenie}[Paraboliczne zagadnienie Dirichleta]\label{thm:pde20-existence}
Niech $\Omega,T,A,f,u_0$ spełniają założenia
definicji~\ref{def:pde20-data}. Wtedy istnieje dokładnie jedno rozwiązanie
słabe zagadnienia~\eqref{eq:pde20-eq} w sensie
definicji~\ref{def:pde20-weak}. Spełnia ono oszacowanie
\begin{equation}\label{eq:pde20-apriori}
 \sup_{0\le t\le T}\|u(t)\|_H^2
 +\lambda\int_0^T\|u(t)\|_V^2\,\mathrm dt
 \le 2\|u_0\|_H^2
 +2\lambda^{-1}\int_0^T\|f(t)\|_{V^*}^2\,\mathrm dt.
\end{equation}
Stała w nierówności nie zależy od gładkości współczynników.
\end{twierdzenie}
\begin{proof}
\emph{Krok 1: równanie skończenie wymiarowe.}
Niech $V_n=\operatorname{span}\{e_1,\ldots,e_n\}$, gdzie $e_j$ są
funkcjami własnymi użytymi w lemacie~\ref{lem:pde20-energy}.
Szukamy $u_n(t)=\sum_{j=1}^n c_j(t)e_j$, które spełnia
\[
 (u_n'(t),e_i)_H+a(t;u_n(t),e_i)
 =\langle f(t),e_i\rangle,\quad i=1,\ldots,n,
 \qquad u_n(0)=P_nu_0.
\]
Jest to liniowy układ $c'(t)+B_n(t)c(t)=F_n(t)$,
gdzie $(B_n(t))_{ij}=a(t;e_j,e_i)$ oraz
$(F_n(t))_i=\langle f(t),e_i\rangle$.
Mierzalność i oszacowania formy z sekcji pierwszej dają
$|(B_n(t))_{ij}|\le\Lambda\sqrt{\mu_i\mu_j}$ oraz
$|(F_n(t))_i|\le\sqrt{\mu_i}\|f(t)\|_{V^*}$.
W konsekwencji $M_n=\operatorname*{ess\,sup}_{t\in(0,T)}
\|B_n(t)\|_{\mathcal L(\R^n)}<\infty$ i
$F_n\in L^2(0,T;\R^n)$.
Na przedziale $[a,b]\subset[0,T]$ o długości
$(b-a)M_n<1$ i dla ustalonego $c_a$ definiujemy
\[
 (\mathcal Tc)(t)=c_a+\int_a^t
 (F_n(r)-B_n(r)c(r))\,\mathrm d\lambda_1(r).
\]
Jest to odwzorowanie zupełnej przestrzeni
$C([a,b];\R^n)$ w siebie: całkowana funkcja należy
do $L^1$, a całka nieoznaczona jest absolutnie ciągła.
Ponadto $\|\mathcal Tc-\mathcal T\widetilde c\|_\infty
\le(b-a)M_n\|c-\widetilde c\|_\infty$.
Twierdzenie Banacha o punkcie stałym daje dokładnie
jedno rozwiązanie równania całkowego.
Dzieląc $[0,T]$ na skończenie wiele takich przedziałów,
za każdym razem bierzemy wartość końcową poprzedniego
rozwiązania jako dane początkowe następnego.
Otrzymujemy jednoznaczne absolutnie ciągłe $c$
na $[0,T]$, z $c_j(0)=(u_0,e_j)_H$.
Ponieważ $c$ jest ograniczona na tym przedziale,
$c'=F_n-B_nc\in L^2(0,T;\R^n)$, czyli $c\in H^1(0,T;\R^n)$.

\emph{Krok 2: energia.}
Mnożymy $i$-te równanie przez $c_i(t)$ i sumujemy po $i$.
Otrzymujemy dla prawie każdego $t$
\[
 \frac12\frac{\mathrm d}{\mathrm dt}\|u_n\|_H^2
 +\lambda\|u_n\|_V^2
 \le\|f\|_{V^*}\|u_n\|_V
 \le\frac{\lambda}{2}\|u_n\|_V^2
 +\frac1{2\lambda}\|f\|_{V^*}^2.
\]
Całkując od $0$ do $t$ i korzystając z
$\|P_nu_0\|_H\le\|u_0\|_H$, dostajemy
\begin{equation}\label{eq:pde20-energy-n}
 \|u_n(t)\|_H^2+\lambda\int_0^t\|u_n(r)\|_V^2\,\mathrm dr
 \le\|u_0\|_H^2+\lambda^{-1}\|f\|_{L^2(0,T;V^*)}^2.
\end{equation}
To jednocześnie zapobiega utracie kontroli nad rozwiązaniami przy $n\to\infty$.

\emph{Krok 3: kontrola pochodnej czasowej.}
Dla każdego $v\in V$ zachodzi
$\langle u_n',v\rangle=\langle f,P_nv\rangle-a(t;u_n,P_nv)$,
ponieważ $u_n'\in V_n$ i $P_n$ jest ortogonalny w $H$.
Korzystamy z $\|P_nv\|_V\le\|v\|_V$ oraz ograniczoności formy:
\[
 \|u_n'(t)\|_{V^*}\le\|f(t)\|_{V^*}+\Lambda\|u_n(t)\|_V.
\]
Zatem $\{u_n'\}$ jest ograniczony w $L^2(0,T;V^*)$.

\emph{Krok 4: granica.}
Po wybraniu podciągu, dzięki refleksywności przestrzeni Hilberta,
$u_n\rightharpoonup u$ w $L^2(0,T;V)$,
$u_n'\rightharpoonup g$ w $L^2(0,T;V^*)$.
Ponadto $u_n$ jest ograniczony w $L^\infty(0,T;H)$.
Dla dowolnego $v\in V$ i $\eta\in C_c^\infty(0,T)$ mamy
\[
 \int_0^T\langle u_n'(t),v\rangle\eta(t)\,\mathrm dt
 =-\int_0^T(u_n(t),v)_H\eta'(t)\,\mathrm dt.
\]
Obie całki są wartościami ciągłych funkcjonałów
liniowych odpowiednio na $L^2(0,T;V^*)$ i $L^2(0,T;V)$.
Przechodząc do granicy dla każdego $v$ i $\eta$,
stwierdzamy, że $g=u'$ w sensie dystrybucyjnym
o wartościach w $V^*$.
Ustalmy $v=e_j$. Dla $n\ge j$ równanie Galerkina
mnożymy przez $\eta\in L^2(0,T)$ i całkujemy po czasie.
Funkcjonał
\[
 w\longmapsto\int_0^T a(t;w(t),e_j)\eta(t)
 \,\mathrm d\lambda_1(t)
\]
jest dobrze określony dzięki mierzalności formy
i ma normę nie większą niż
$\Lambda\|e_j\|_V\|\eta\|_{L^2(0,T)}$.
Słaba zbieżność daje więc
\[
 \int_0^T
 (\langle u',e_j\rangle+a(t;u,e_j)-\langle f,e_j\rangle)
 \eta\,\mathrm d\lambda_1=0.
\]
Funkcja w nawiasie należy do $L^2(0,T)$; wybór jej
samej jako $\eta$ dowodzi, że znika prawie wszędzie.
Po odrzuceniu przeliczalnej sumy zbiorów wyjątkowych
równość zachodzi jednocześnie dla wszystkich $e_j$.
Przez liniowość, gęstość ich rozpiętości w $V$
i ciągłość funkcjonałów otrzymujemy równanie w $V^*$
dla prawie każdego $t$.

\emph{Krok 5: warunek początkowy.}
Dla $v\in V_n$ mnożymy równanie Galerkina przez
dowolną $\eta\in C^1([0,T])$ spełniającą $\eta(T)=0$,
całkujemy od $0$ do $T$ i całkujemy pierwszy składnik
przez części. Otrzymujemy
\[
 -\int_0^T(u_n(t),v)_H\eta'(t)\,\mathrm dt
 +\int_0^T a(t;u_n(t),v)\eta(t)\,\mathrm dt
 =\int_0^T\langle f(t),v\rangle\eta(t)\,\mathrm dt
 +(P_nu_0,v)_H\eta(0).
\]
Dla każdego ustalonego $v=e_j$ przechodzimy do granicy jak wyżej.
Z lematu~\ref{lem:pde20-energy} funkcja
$t\mapsto(u(t),e_j)_H$ ma absolutnie ciągły reprezentant
o pochodnej $\langle u'(t),e_j\rangle$.
Całkowanie równania granicznego przez części
dla tego współczynnika daje tę samą tożsamość
z wyrazem $(u(0),e_j)_H\eta(0)$ zamiast
$(u_0,e_j)_H\eta(0)$.
Odejmując obie tożsamości i biorąc np.
$\eta(t)=1-t/T$, otrzymujemy
$(u(0),e_j)_H=(u_0,e_j)_H$ dla każdego $j$.
Zatem $u(0)=u_0$ w $H$.

\emph{Krok 6: oszacowanie i jednoznaczność.}
Wstawiamy $v=u(t)$ do równania słabego (jest to dopuszczalne
w równości w $V^*$ dla prawie każdego $t$) i używamy
lematu~\ref{lem:pde20-energy}. Z koercywności i nierówności
Younga, tak jak w kroku 2, dla każdego $t$ dostajemy
\[
 \|u(t)\|_H^2+\lambda\int_0^t\|u(r)\|_V^2\,\mathrm d\lambda_1(r)
 \le M,\qquad
 M=\|u_0\|_H^2+\lambda^{-1}\|f\|_{L^2(0,T;V^*)}^2.
\]
Pierwszy składnik jest ograniczony przez $M$ jednostajnie
w czasie, a drugi przez $M$ dla $t=T$. Dodanie tych
dwóch oszacowań daje~\eqref{eq:pde20-apriori}.
Z równania w $V^*$ otrzymujemy również
$\|u'\|_{L^2(0,T;V^*)}\le
\|f\|_{L^2(0,T;V^*)}+\Lambda\|u\|_{L^2(0,T;V)}$.
Jeżeli $u,w$ są rozwiązaniami z tymi samymi danymi,
to $z=u-w$ spełnia równanie z zerową prawą stroną
i zerową wartością początkową. Tożsamość energii daje
$\|z(t)\|_H^2+2\lambda\int_0^t\|z\|_V^2\,\mathrm dr\le0$.
Stąd $z=0$.
\end{proof}

\begin{uwaga}[Zmienne w czasie współczynniki i brak symetrii]
W dowodzie nie różniczkowaliśmy $A$ względem czasu
ani nie odwoływaliśmy się do symetrii $a(t;\cdot,\cdot)$.
Jest to istotna przewaga słabego sformułowania nad próbami
obliczania $\Delta u$ punkt po punkcie.
\end{uwaga}

\section{Zależność od danych i paraboliczna zasada maksimum}
\begin{wniosek}[Stabilność rozwiązania]\label{cor:pde20-stability}
Jeżeli $u_i$ odpowiadają danym $(u_{0,i},f_i)$, $i=1,2$,
przy tej samej macierzy $A$, to dla $z=u_1-u_2$ mamy
\[
 \sup_{t\in[0,T]}\|z(t)\|_H^2
 +\lambda\|z\|_{L^2(0,T;V)}^2
 \le2\|u_{0,1}-u_{0,2}\|_H^2
 +2\lambda^{-1}\|f_1-f_2\|_{L^2(0,T;V^*)}^2.
\]
\end{wniosek}
\begin{proof}
Równanie na $z$ ma tę samą formę i odpowiadającą jej prawą
stronę $f_1-f_2$. Stosujemy twierdzenie~\ref{thm:pde20-existence}.
\end{proof}

\begin{lemat}[Reguła obcięcia w czasie]\label{lem:pde20-trunc}
Jeżeli $u\in L^2(0,T;V)$ i $u'\in L^2(0,T;V^*)$, to
$u^-:=\max\{-u,0\}\in L^2(0,T;V)$ i funkcja
$t\mapsto\|u^-(t)\|_H^2$ jest absolutnie ciągła, przy czym
\[
 \frac12\frac{\mathrm d}{\mathrm dt}\|u^-(t)\|_H^2
 =-\langle u'(t),u^-(t)\rangle
 \quad\text{dla prawie wszystkich }t.
\]
\end{lemat}
\begin{proof}
Niech $N(v)=v^-$. Z reguł obcięcia
~\ref{cor:pde9-truncation} oraz zanikania gradientu
na zbiorze poziomicowym~\ref{lem:pde9-levelsets}
mamy $\nabla N(v)=-\chi_{\{v<0\}}\nabla v$
prawie wszędzie. Dla $v\in V$ obcięcie należy do $V$:
przybliżamy $v$ funkcjami z $C_c^\infty(\Omega)$,
bierzemy ich części ujemne (funkcje Sobolewa o zwartym
nośniku, należące do $V$ przez wygładzanie splotem).
Są ograniczone w $V$ i zbiegają do $N(v)$ w $H$,
ponieważ $N$ jest lipschitzowskie ze stałą $1$.
Słabo zbieżny podciąg w $V$ ma więc granicę $N(v)$.
Ponadto $\|N(v)\|_V\le\|v\|_V$.

Jeśli $v_n\to v$ w $L^2(0,T;V)$, to
$N(v_n)\to N(v)$ w tej samej przestrzeni. Istotnie,
\[
 \begin{split}
 \nabla N(v_n)-\nabla N(v)
 &=-\chi_{\{v_n<0\}}(\nabla v_n-\nabla v)\\
 &\quad+(\chi_{\{v<0\}}-\chi_{\{v_n<0\}})\nabla v.
 \end{split}
\]
Pierwszy składnik dąży do zera w $L^2$ na
$(0,T)\times\Omega$. Z każdego podciągu wybieramy
dalszy podciąg, dla którego $v_n\to v$ prawie wszędzie.
Poza $\{v=0\}$ funkcje charakterystyczne w drugim
składniku zbiegają, a na $\{v=0\}$ gradient $v$
znika prawie wszędzie (stosujemy lemat o poziomicach
dla prawie każdego czasu i twierdzenie Fubiniego).
Z twierdzenia o zbieżności zmajoryzowanej, z majorantą
$|\nabla v|^2$ dla kwadratu modułu, drugi składnik
dąży do zera w $L^2$ wzdłuż tego dalszego podciągu.
Argument podciągowy daje zbieżność całego ciągu.
Analogicznie dla pojedynczej funkcji czasu otrzymujemy
ciągłość $N:V\to V$, a stąd silną mierzalność $N(u)$
oraz $N(u)\in L^2(0,T;V)$.

Na $H$ określamy funkcjonał
$J(v)=\frac12\|v^-\|_H^2$.
Funkcja skalarna $\Phi(s)=\frac12(s^-)^2$ jest klasy
$C^1$ i $\Phi'(s)=-s^-$ jest lipschitzowska
ze stałą $1$. Z twierdzenia Newtona--Leibniza wynika
$|\Phi(a+b)-\Phi(a)-\Phi'(a)b|\le\frac12|b|^2$.
Całkując, dostajemy
\[
 |J(v+h)-J(v)+(v^-,h)_H|\le\tfrac12\|h\|_H^2.
\]
Zatem $J$ jest różniczkowalny w sensie Frécheta,
$DJ(v)h=-(v^-,h)_H$, a jego różniczka jest ciągła
dzięki lipschitzowskości $N$ w $H$.

Stosujemy teraz rzuty $P_n$ z dowodu lematu
~\ref{lem:pde20-energy}. Funkcje $u_n=P_nu$ są
absolutnie ciągłe o wartościach w skończenie
wymiarowej przestrzeni $V_n$ i $u_n'=P_nu'$.
Reguła łańcuchowa dla funkcji absolutnie ciągłych
i funkcjonału $J$ daje dla $0\le s\le t\le T$
\[
 J(u_n(t))-J(u_n(s))
 =-\int_s^t\langle P_nu'(r),N(u_n(r))\rangle
 \,\mathrm d\lambda_1(r).
\]
Iloczyn skalarny w $H$ jest tu równy dualności,
ponieważ $u_n'\in V_n$ prawie wszędzie.
Z lematu~\ref{lem:pde20-energy} mamy zbieżność
$u_n\to u$ w $C([0,T];H)$ i $L^2(0,T;V)$
oraz $P_nu'\to u'$ w $L^2(0,T;V^*)$.
Wykazana ciągłość obcięcia daje
$N(u_n)\to N(u)$ w $L^2(0,T;V)$.
Stąd przez nierówność Cauchy'ego--Schwarza
całki w dualności zbiegają do całki
$-\int_s^t\langle u',N(u)\rangle\,\mathrm d\lambda_1$,
a wartości końcowe $J(u_n(s)),J(u_n(t))$ zbiegają
dzięki ciągłości $J$ na $H$.
Otrzymujemy tożsamość całkową dla wszystkich $s,t$.
Ponieważ $\langle u',N(u)\rangle\in L^1(0,T)$,
wynika z niej absolutna ciągłość $J(u(\cdot))$
i żądana równość po zróżniczkowaniu prawie wszędzie.
\end{proof}

\begin{twierdzenie}[Zachowanie dodatniości]\label{thm:pde20-positivity}
Niech $u$ będzie rozwiązaniem z twierdzenia~\ref{thm:pde20-existence}.
Jeżeli $u_0\ge0$ prawie wszędzie i $\langle f(t),v\rangle\ge0$
dla każdej nieujemnej $v\in V$ i prawie każdego $t$, to
$u(t,x)\ge0$ dla prawie każdego $x\in\Omega$ i każdego $t\in[0,T]$.
\end{twierdzenie}
\begin{proof}
Stożek nieujemnych elementów $V$ ma przeliczalny
gęsty podzbiór, ponieważ $V$ jest ośrodkową
przestrzenią metryczną. Odrzucając sumę zbiorów
wyjątkowych dla tego podzbioru i korzystając
z ciągłości $f(t)$ na $V$, możemy przyjąć, że
$\langle f(t),v\rangle\ge0$ dla wszystkich
nieujemnych $v\in V$ poza jednym zbiorem miary zero.
Testujemy równanie w $V^*$ przez $u^-(t)$
dla prawie każdego $t$. Z lematu~\ref{lem:pde20-trunc}
$\langle u',u^-\rangle=-\frac12\mathrm d\|u^-\|_H^2/\mathrm dt$.
Na zbiorze $\{u<0\}$ mamy $\nabla u=-\nabla u^-$,
a na dopełnieniu $\nabla u^-=0$ prawie wszędzie. Zatem
$a(t;u,u^-)=-a(t;u^-,u^-)\le-\lambda\|u^-\|_V^2$.
W rezultacie
\[
 \frac12\frac{\mathrm d}{\mathrm dt}\|u^-(t)\|_H^2
 +\lambda\|u^-(t)\|_V^2\le0.
\]
Ponieważ $u^-(0)=0$, lewa strona po scałkowaniu daje
$\|u^-(t)\|_H^2=0$ dla każdego $t$.
\end{proof}

\begin{wniosek}[Zasada porównania]
Dla tej samej macierzy $A$ niech rozwiązania $u,v$ mają
odpowiednio dane $(u_0,f)$ i $(v_0,g)$.
Jeżeli $u_0\le v_0$ prawie wszędzie i
$\langle g(t)-f(t),\phi\rangle\ge0$ dla wszystkich
nieujemnych $\phi\in V$ i prawie wszystkich $t$,
to $u(t)\le v(t)$ prawie wszędzie dla każdego $t$.
\end{wniosek}
\begin{proof}
Stosujemy twierdzenie~\ref{thm:pde20-positivity}
do różnicy $v-u$.
\end{proof}

\section{Wygładzanie i zanik energii dla laplasjanu Dirichleta}
Niech $L=-\Delta_D$ będzie dodatnim laplasjanem Dirichleta
w $H$. Jego wartości własne oznaczamy przez
$0<\mu_1\le\mu_2\le\cdots\to\infty$ (z krotnościami),
a ortonormalne funkcje własne przez $e_j$
(twierdzenie~\ref{thm:pde18-compact}).
Dla $s\ge0$ określamy potęgę spektralną przez
\[
 D(L^s)=\left\{v\in H:\sum_{j\ge1}\mu_j^{2s}|(v,e_j)_H|^2<\infty\right\},
 \qquad L^sv=\sum_{j\ge1}\mu_j^s(v,e_j)_He_j.
\]
W szczególności $D(L^{1/2})=V$ i
$\|L^{1/2}v\|_H=\|v\|_V$ z rozwinięcia formy
w dowodzie twierdzenia~\ref{thm:pde18-minmax}.
Odróżniamy operator $L:D(L)\subset H\to H$
od jego rozszerzenia formowego $\mathcal L:V\to V^*$,
określonego przez
$\langle\mathcal Lv,w\rangle=(\nabla v,\nabla w)_{L^2}$.
Mamy $\|\mathcal Lv\|_{V^*}=\|v\|_V$, co wynika
z twierdzenia Riesza w $V$.
W tej sekcji zakładamy
$A=I$ i $f=0$, chyba że wyraźnie zaznaczymy inaczej.

\begin{twierdzenie}[Półgrupa cieplna i wygładzanie]\label{thm:pde20-semigroup}
Dla każdego $u_0\in H$ funkcja
$u(t)=e^{-tL}u_0=\sum_{j\ge1}e^{-\mu_jt}(u_0,e_j)_He_j$
jest jedynym rozwiązaniem słabym zagadnienia jednorodnego.
Dla każdego $s\ge0$ i $t>0$ zachodzi
\[
 \|L^se^{-tL}u_0\|_H
 \le C_s t^{-s}\|u_0\|_H,
 \qquad C_s=(s/\mathrm e)^s\ (s>0),\quad C_0=1.
\]
Ponadto $\|u(t)\|_H\le e^{-\mu_1t}\|u_0\|_H$.
\end{twierdzenie}
\begin{proof}
Dla skończonych sum własnych obliczamy
$u'(t)+Lu(t)=0$, a tożsamość energii jest bezpośrednia.
Niech $b_j=(u_0,e_j)_H$. Dla ogólnego $u_0$
kwadrat normy ogona szeregu w $H$ nie przewyższa
$\sum_{j>N}|b_j|^2$, jednostajnie dla $t\ge0$.
Szereg określa więc ciągłą funkcję na $[0,\infty)$,
przy czym $u(0)=u_0$.
Mnożenie współczynników przez $e^{-\mu_jt}$ określa
ograniczony operator $S(t)$ na $H$. Tożsamość
$e^{-\mu_j(t+r)}=e^{-\mu_jt}e^{-\mu_jr}$ daje
$S(t+r)=S(t)S(r)$ oraz $S(0)=I$.
Jest to silnie ciągła półgrupa, a zapis $e^{-tL}$
oznacza właśnie tę konstrukcję spektralną.
Z twierdzenia Tonellego dla nieujemnych składników mamy
\[
 \int_0^T\|L^{1/2}e^{-tL}u_0\|_H^2\,\mathrm dt
 =\sum_j\mu_j|(u_0,e_j)_H|^2
 \int_0^T e^{-2\mu_jt}\,\mathrm dt
 \le\tfrac12\|u_0\|_H^2.
\]
To samo oszacowanie dla ogonów pokazuje, że
skończone sumy $u_N$ zbiegają w $L^2(0,T;V)$.
Ich pochodne $u_N'=-\mathcal Lu_N$ zbiegają
w $L^2(0,T;V^*)$ do $-\mathcal Lu$ dzięki izometrii
$\mathcal L$. Przechodząc do granicy w definicji
pochodnej dystrybucyjnej, otrzymujemy
$u'=-\mathcal Lu$. Razem z warunkiem początkowym
dowodzi to, że granica spełnia definicję rozwiązania słabego;
jednoznaczność wynika z twierdzenia~\ref{thm:pde20-existence}.
Dla $s>0$ supremum funkcji $y\mapsto y^se^{-ty}$ na
$[0,\infty)$ wynosi $(s/(\mathrm et))^s$: pochodna
dla $y>0$ ma znak $s-ty$, a funkcja dąży do zera
na obu krańcach. Dla $s=0$ używamy oszacowania
$e^{-ty}\le1$. Zatem
\[
 \sum_j\mu_j^{2s}e^{-2\mu_jt}|b_j|^2
 \le C_s^2t^{-2s}\sum_j|b_j|^2.
\]
Dowodzi to jednocześnie $u(t)\in D(L^s)$ i nierówności
wygładzającej. Ponieważ $\mu_j\ge\mu_1$,
z tożsamości Parsevala otrzymujemy
$\|u(t)\|_H^2\le e^{-2\mu_1t}\|u_0\|_H^2$.
Ponadto dla $0\le r\le t$ dokładne całkowanie
każdego składnika i twierdzenie Tonellego dają
\[
 \|u(t)\|_H^2+2\int_r^t\|u(\sigma)\|_V^2
 \,\mathrm d\lambda_1(\sigma)=\|u(r)\|_H^2.
\]
\end{proof}


\begin{twierdzenie}[Wzór Duhamela dla wymuszenia w $L^2(H)$]
\label{thm:pde20-duhamel}
Niech $\Omega\subset\R^d$ będzie ograniczonym obszarem
Lipschitzowskim, $H=L^2(\Omega)$, $L=-\Delta_D$,
$u_0\in H$, $f\in L^2(0,T;H)$.
Jedynym rozwiązaniem słabym zagadnienia $u'+Lu=f$,
$u(0)=u_0$, jest funkcja
\[
 u(t)=e^{-tL}u_0+\int_0^t e^{-(t-s)L}f(s)\,\mathrm ds,
 \qquad 0\le t\le T,
\]
gdzie całka jest całką Bochnera w $H$.
\end{twierdzenie}
\begin{proof}
Operator $S(t)=e^{-tL}$ jest kontrakcją w $H$.
Dla każdego ustalonego $t\in[0,T]$ funkcja
$s\mapsto S(t-s)f(s)$ na $[0,t]$ jest silnie
mierzalna: dla funkcji prostych $f$ wynika to
z silnej ciągłości $S$, a dla ogólnej $f$
z aproksymacji punktowej prawie wszędzie.
Jej norma jest ograniczona przez $\|f(s)\|_H$,
która należy do $L^1(0,T)$, gdyż
$\|f\|_{L^1(0,T;H)}\le\sqrt T\|f\|_{L^2(0,T;H)}$.
Całka Bochnera istnieje więc dla każdego $t$.
Jeżeli $t_n\to t$, przedłużamy integrandy zerem
na $(t_n,T)$; dla prawie każdego $s\ne t$
zbiegają silnie w $H$ do odpowiedniego integrandu
dla $t$. Twierdzenie o zbieżności zmajoryzowanej
dla całek Bochnera daje ciągłość składnika całkowego.

Niech $u_n$ będzie rozwiązaniem w $V_n$ z danymi
$P_nu_0$ i $f_n=P_nf$. Dla $1\le j\le n$
jego współczynniki spełniają
$r_j'+\mu_jr_j=(f,e_j)_H$,
$r_j(0)=(u_0,e_j)_H$. Mnożymy przez $e^{\mu_jt}$
i całkujemy, uzyskując
$r_j(t)=e^{-\mu_jt}b_j+
\int_0^t e^{-\mu_j(t-s)}(f(s),e_j)_H
\,\mathrm d\lambda_1(s)$.
Sumując skończenie wiele współczynników, dostajemy
wzór Duhamela z danymi $P_nu_0,P_nf$. Dla różnic dwóch skończonych przybliżeń
stosujemy oszacowanie energii z
Twierdzenia~\ref{thm:pde20-existence}:
\[
 \sup_{0\le t\le T}\|u_n(t)-u_m(t)\|_H^2
 +\int_0^T\|u_n-u_m\|_V^2\,\mathrm dt
 \le C\bigl(\|P_nu_0-P_mu_0\|_H^2
 +\|P_nf-P_mf\|_{L^2(0,T;H)}^2\bigr),
\]
przy czym w prawej stronie korzystamy z ciągłego
zanurzenia $H\hookrightarrow V^*$. Prawa strona dąży do
zera. Równanie dla różnic daje również
$\|u_n'-u_m'\|_{L^2(0,T;V^*)}\le
 C_P\|P_nf-P_mf\|_{L^2(0,T;H)}
 +\|u_n-u_m\|_{L^2(0,T;V)}$.
Z zupełności i przejścia granicznego w definicji
pochodnej wynika, że granica spełnia równanie
$u'+\mathcal Lu=f$ w $L^2(0,T;V^*)$.
Warunek początkowy przechodzi dzięki zbieżności
w $C([0,T];H)$. Jednocześnie różnica między wzorem
całkowym dla $u_n$ a wzorem z tezy ma normę w
$C([0,T];H)$ nie większą niż
$\|P_nu_0-u_0\|_H+\sqrt T\|P_nf-f\|_{L^2(0,T;H)}$.
Obie granice są więc identyczne; jednoznaczność
daje twierdzenie~\ref{thm:pde20-existence}.
\end{proof}

\begin{twierdzenie}[Wyższa regularność energii]\label{thm:pde20-max-reg}
Niech $L=-\Delta_D$, $u_0\in V=D(L^{1/2})$,
$f\in L^2(0,T;H)$. Wtedy rozwiązanie słabe równania
$u'+Lu=f$ spełnia
$u\in H^1(0,T;H)\cap L^2(0,T;D(L))$ oraz
\[
 \int_0^T\|u'(t)\|_H^2\,\mathrm dt
 +\sup_{t\in[0,T]}\|L^{1/2}u(t)\|_H^2
 \le2\|L^{1/2}u_0\|_H^2
 +2\int_0^T\|f(t)\|_H^2\,\mathrm dt.
\]
Dla obszarów klasy $C^{1,1}$ można ponadto utożsamić
$D(L)=H^2(\Omega)\cap H_0^1(\Omega)$; bez regularności brzegu
takiego utożsamienia nie zakładamy.
\end{twierdzenie}
\begin{proof}
Dla $u_n=P_nu$ (równoważnie skończonych rozwiązań
z dowodu twierdzenia~\ref{thm:pde20-duhamel})
korzystamy z równania $u_n'+Lu_n=P_nf$
z początkową daną $P_nu_0$.
Iloczyn skalarny z $u_n'$ daje
$\|u_n'\|_H^2+\frac12\frac{\mathrm d}{\mathrm dt}
\|L^{1/2}u_n\|_H^2=(P_nf,u_n')_H$.
Stosujemy nierówność
$|(P_nf,u_n')|\le(\|f\|_H^2+\|u_n'\|_H^2)/2$,
po czym całkujemy:
\[
 \int_0^t\|u_n'(r)\|_H^2\,\mathrm d\lambda_1(r)
 +\|L^{1/2}u_n(t)\|_H^2
 \le \|P_nu_0\|_V^2+\int_0^t\|P_nf(r)\|_H^2\,\mathrm d\lambda_1(r).
\]
To samo oszacowanie stosujemy do $u_n-u_m$, z danymi
$(P_n-P_m)u_0$ i $(P_n-P_m)f$. Prawa strona
dąży do zera z silnej zbieżności rzutów w $V$
oraz $L^2(0,T;H)$. Otrzymujemy zatem ciąg Cauchy'ego
w $C([0,T];V)$ oraz ciąg pochodnych Cauchy'ego
w $L^2(0,T;H)$. Równanie pokazuje ponadto, że
$Lu_n$ jest ciągiem Cauchy'ego w $L^2(0,T;H)$.
Przy normie grafowej
$\|v\|_{D(L)}^2=\|v\|_H^2+\|Lv\|_H^2$
przestrzeń $D(L)$ jest zupełna dzięki domkniętości $L$.
W konsekwencji $u_n$ zbiega również w
$L^2(0,T;D(L))$. Wszystkie granice funkcji $u_n$
są równe $u$, gdyż zanurzenia do $L^2(0,T;H)$
są ciągłe. Przechodząc do granicy w tożsamości
dystrybucyjnej dla pochodnych, otrzymujemy
$u\in H^1(0,T;H)$ i $Lu=f-u'$ w $L^2(0,T;H)$.
Mamy także $u\in C([0,T];V)$.
Oznaczmy $M=\|u_0\|_V^2+\|f\|_{L^2(0,T;H)}^2$.
Przejście do granicy w oszacowaniu daje
$\|u'\|_{L^2(0,T;H)}^2\le M$ oraz
$\sup_t\|u(t)\|_V^2\le M$; dodanie obu nierówności
daje tezę. Dodatkowo
$\|Lu\|_{L^2(0,T;H)}\le\|f\|_{L^2(0,T;H)}
+\|u'\|_{L^2(0,T;H)}$.
\end{proof}

\section{Zwartość czasowo-przestrzenna: lemat Aubina--Lionsa w naszym modelu}
\begin{twierdzenie}[Zwartość dla parabolicznych ciągów Galerkina]\label{thm:pde20-compact}
Niech $\Omega\subset\R^d$ będzie ograniczonym obszarem
Lipschitzowskim. Jeżeli ciąg $\{u_n\}$ jest ograniczony
w $L^2(0,T;V)$ i ciąg $\{u_n'\}$ jest ograniczony
w $L^2(0,T;V^*)$, to $\{u_n\}$ jest względnie zwarty
w $L^2(0,T;H)$.
\end{twierdzenie}
\begin{proof}
Niech $P_N$ oznacza stały rzut na pierwsze $N$ funkcji
własnych laplasjanu Dirichleta. Z rozwinięcia w bazie
$\{e_j\}$ dostajemy dla prawie każdego $t$
\[
 \|(I-P_N)u_n(t)\|_H^2
 \le\mu_{N+1}^{-1}\|u_n(t)\|_V^2.
\]
Całkowanie po czasie i wspólne ograniczenie norm $V$ dają
\begin{equation}\label{eq:pde20-tail}
 \sup_n\|(I-P_N)u_n\|_{L^2(0,T;H)}^2
 \le C\mu_{N+1}^{-1}\xrightarrow[N\to\infty]{}0.
\end{equation}
Dla ustalonego $N$ współczynniki
$c_{n,j}(t)=(u_n(t),e_j)_H$, $j\le N$, są ograniczone
w $L^2(0,T)$, a ich pochodne spełniają
$|c_{n,j}'(t)|=|\langle u_n'(t),e_j\rangle|
\le\|u_n'(t)\|_{V^*}\|e_j\|_V$.
Zatem dla każdego $j\le N$ ciąg $\{c_{n,j}\}_n$ jest
ograniczony w $H^1(0,T)$, a więc względnie zwarty
w $C([0,T])$. Dla absolutnie ciągłego reprezentanta
nierówność Cauchy'ego--Schwarza daje
$|c(t)-c(s)|\le\|c'\|_{L^2}|t-s|^{1/2}$.
Wybierając $s$ z $|c(s)|\le T^{-1/2}\|c\|_{L^2}$,
otrzymujemy również
$\|c\|_\infty\le T^{-1/2}\|c\|_{L^2}
+\sqrt T\|c'\|_{L^2}$.
Mamy zatem jednostajne ograniczenie i równociągłość,
więc stosujemy twierdzenie Arzelà--Ascoliego.
Dla każdego $N$ wybieramy kolejno podciąg poprzedniego
podciągu, na którym pierwsze $N$ współczynników
zbiegają w $C([0,T])$. Podciąg przekątny $(u_{n_k})$
ma tę własność dla każdego ustalonego $N$,
a zatem $P_Nu_{n_k}$ zbiega w $L^2(0,T;H)$.
Dla $\eps>0$ wybieramy $N$ tak, by z~\eqref{eq:pde20-tail}
$\sup_n\|(I-P_N)u_n\|_{L^2(0,T;H)}<\eps/3$.
Następnie dla dostatecznie dużych $k,l$
$\|P_N(u_{n_k}-u_{n_l})\|_{L^2(0,T;H)}<\eps/3$.
Nierówność trójkąta daje
$\|u_{n_k}-u_{n_l}\|_{L^2(0,T;H)}<\eps$.
Z zupełności tej przestrzeni podciąg jest zbieżny,
co dowodzi względnej zwartości.
\end{proof}


\begin{wniosek}[Szacunek przesunięć czasowych w normie $H$]
\label{cor:pde20-translations}
Niech $\Omega\subset\R^d$ będzie ograniczonym obszarem
Lipschitzowskim, $V=H_0^1(\Omega)$, $H=L^2(\Omega)$,
a $u\in L^2(0,T;V)$, $u'\in L^2(0,T;V^*)$.
Dla $0<h<T$ zachodzi
\[
 \int_0^{T-h}\|u(t+h)-u(t)\|_H^2\,\mathrm dt
 \le 2h\,\|u\|_{L^2(0,T;V)}
 \|u'\|_{L^2(0,T;V^*)}.
\]
W szczególności rodzina funkcji o jednostajnie
ograniczonych obu normach ma jednostajny moduł ciągłości
względem translacji czasu w $L^2(H)$.
\end{wniosek}
\begin{proof}
Używamy ciągłego reprezentanta z lematu
~\ref{lem:pde20-energy}. Dla skończonych rzutów
mamy tożsamość
$P_Nu(t+h)-P_Nu(t)=\int_t^{t+h}P_Nu'(s)
\,\mathrm d\lambda_1(s)$ w $V^*$.
Przechodzimy do granicy: lewa strona zbiega dzięki
jednostajnej zbieżności w $H$ i ciągłemu zanurzeniu
do $V^*$, a prawa z nierówności
$\|\int_t^{t+h}(P_Nu'-u')\,\mathrm d\lambda_1\|_{V^*}
\le\sqrt h\|P_Nu'-u'\|_{L^2(0,T;V^*)}$.
Dla wszystkich $t\in[0,T-h]$ otrzymujemy w $V^*$
$u(t+h)-u(t)=\int_t^{t+h}u'(s)\,\mathrm d\lambda_1(s)$.
Dzięki nierówności Minkowskiego dla przesunięć oraz
przedłużeniu $u'$ przez zero poza $(0,T)$ mamy
\[
 \|u(\cdot+h)-u\|_{L^2(0,T-h;V^*)}
 \le h\|u'\|_{L^2(0,T;V^*)}.
\]
Z drugiej strony nierówność trójkąta daje
$\|u(\cdot+h)-u\|_{L^2(0,T-h;V)}
\le2\|u\|_{L^2(0,T;V)}$.
Z osadzeń $V\hookrightarrow H\hookrightarrow V^*$,
przy identyfikacji $H=H^*$, otrzymujemy
$\|w\|_H^2\le\|w\|_V\|w\|_{V^*}$.
Całkujemy tę nierówność dla $w(t)=u(t+h)-u(t)$,
a następnie stosujemy Cauchy'ego--Schwarza w czasie.
\end{proof}

\begin{uwaga}[Zastosowanie do nieliniowości]
Zbieżność słaba w $L^2(0,T;V)$ na ogół nie pozwala
przechodzić do granicy w wyrażeniu nieliniowym $F(u_n)$.
Zbieżność silna w $L^2(0,T;H)$ otrzymana powyżej
jest często brakującym elementem dowodów istnienia
w nieliniowych PDE. Ogólne twierdzenie Aubina--Lionsa
obejmuje więcej par przestrzeni i inne wykładniki;
tutaj udowodniliśmy samodzielnie wystarczający przypadek.
\end{uwaga}

\section{Nieliniowy kontrast: możliwość eksplozji}
\begin{twierdzenie}[Warunkowa eksplozja dla równania reakcji--dyfuzji]
\label{thm:pde20-blowup}
Niech $p>1$, $\Omega\subset\R^d$ będzie ograniczonym,
gładkim obszarem, a $\phi_1>0$ wewnątrz $\Omega$
będzie pierwszą funkcją własną laplasjanu Dirichleta,
$-\Delta\phi_1=\mu_1\phi_1$, z normalizacją
$\int_\Omega\phi_1\,\mathrm d\lambda_d(x)=1$.
Niech $u\ge0$ będzie dostatecznie regularnym rozwiązaniem
$u_t-\Delta u=u^p$ na $[0,T_*)\times\Omega$
z jednorodnym warunkiem Dirichleta. Jeżeli
$F(0):=\int_\Omega u(0,x)\phi_1(x)\,\mathrm d\lambda_d(x)
>\mu_1^{1/(p-1)}$, to nie istnieje globalne, nieujemne rozwiązanie klasyczne
na $[0,\infty)$ z funkcjonałem $F(t)$ skończonym w każdej chwili.
Dokładniej, jeśli $T_y<\infty$ jest czasem eksplozji rozwiązania
$y'=-\mu_1y+y^p$, $y(0)=F(0)$, to rozwiązanie $u$ nie może
pozostawać klasyczne na żadnym przedziale $[0,T]$ z $T\ge T_y$.
Jeśli istnieje aż do $T_y$, wówczas $F(t)\to\infty$ przy $t\uparrow T_y$.
\end{twierdzenie}
\begin{proof}
Mnożąc równanie przez $\phi_1$ i całkując po $\Omega$,
a następnie korzystając z całkowania przez części
i z warunków Dirichleta, otrzymujemy
\[
 F'(t)=-\mu_1F(t)+\int_\Omega u(t,x)^p\phi_1(x)
 \,\mathrm d\lambda_d(x).
\]
Miara $\phi_1(x)\,\mathrm d\lambda_d(x)$ jest probabilistyczna,
więc nierówność Jensena daje
$\int_\Omega u^p\phi_1\,\mathrm d\lambda_d\ge F(t)^p$.
Stąd $F'\ge-\mu_1F+F^p$.
Dla $F(0)>\mu_1^{1/(p-1)}$ rozwiązanie skalarnego ODE
$y'=-\mu_1y+y^p$ rośnie i osiąga nieskończoność
w skończonym czasie, ponieważ dla dużych $y$
mamy $y^p-\mu_1y\ge y^p/2$, a
$\int_{y_0}^\infty 2s^{-p}\,\mathrm ds<\infty$.
Z zasady porównania dla skalarnych ODE
$F(t)\ge y(t)$, dopóki oba wyrażenia są skończone.
To wyklucza globalne ograniczone rozwiązanie wskazanego typu.
\end{proof}

\section{Stabilność wstecznego Eulera i rozdzielenie czasu od przestrzeni}
Dla dyskretyzacji laplasjanu Dirichleta na siatce jednowymiarowej
niech $L_h$ będzie symetryczną dodatnio określoną macierzą
różnicową. Przy kroku czasowym $\tau>0$ schemat
wstecznego Eulera ma postać
\begin{equation}\label{eq:pde20-implicit}
 \frac{U^{m+1}-U^m}{\tau}+L_hU^{m+1}=F^{m+1}.
\end{equation}
\begin{stwierdzenie}[Bezwarunkowa stabilność energetyczna]
\label{prop:pde20-be}
Jeżeli $F^{m+1}=0$ dla wszystkich $m$, to rozwiązanie
schematu~\eqref{eq:pde20-implicit} spełnia
$\|U^{m+1}\|_h\le\|U^m\|_h$ niezależnie od $\tau/h^2$,
gdzie $\|U\|_h^2=h\sum_j|U_j|^2$.
\end{stwierdzenie}
\begin{proof}
Bierzemy dyskretny iloczyn skalarny równania przez $U^{m+1}$.
Tożsamość $2(a-b,a)=\|a\|^2-\|b\|^2+\|a-b\|^2$
daje
\[
 \|U^{m+1}\|_h^2-\|U^m\|_h^2
 +\|U^{m+1}-U^m\|_h^2
 +2\tau(L_hU^{m+1},U^{m+1})_h=0.
\]
Ostatnie dwa składniki są nieujemne; stąd teza.
Jest to twierdzenie o stabilności, nie samoistny dowód
zbieżności: potrzebna jest jeszcze kontrola błędu aproksymacji.
\end{proof}

\begin{uwaga}[Granice bieżącego rozdziału]
Dla współczynników zmiennych w czasie udowodniliśmy
istnienie i jednoznaczność rozwiązania w naturalnej klasie
$L^2(V)\cap H^1(V^*)$ bez gładkości $A$. Nie dowodziliśmy
tu pełnej parabolicznej teorii De Giorgiego--Nasha--Mosera,
teorii oszacowań Schaudera ani maksymalnej regularności
w klasach $L^p$ dla dowolnych operatorów niesymetrycznych.
Odpowiednie zagadnienia wymagają odrębnych założeń i metod.
\end{uwaga}

